\section{Method}
\label{sec:methods}
Our method addresses the core challenge of uncertainty-aware human-AI collaboration through a principled four-stage framework. Given the increasingly critical role of LLMs in high-stakes content moderation, we propose a supervised approach to uncertainty quantification that learns when to trust model outputs versus when to escalate for human review. This design directly addresses our research questions: \textbf{RQ1} (can we accurately predict LLM errors?) and \textbf{RQ2} (can selective escalation based on these predictions reduce the overall cost in human-AI moderation workflows?).
As illustrated in Figure~\ref{fig:pipeline}, our framework proceeds in four stages:
(1) \firstmention{Base LLM Inference}: multimodal content is processed by a base LLM using structured prompts that enforce deterministic, token-aligned outputs; (2) \firstmention{Integer-Token Output Schema}: to ensure consistent probability extraction across diverse base LLMs, we designed prompts that constrain outputs to single integer tokens: \texttt{0}~=~``no'', \texttt{1}~=~``yes'', \texttt{2}~=~``inconclusive\_evidence'' (aleatoric uncertainty), \texttt{and 3}~=~``inconclusive\_definition'' (epistemic uncertainty), which are validated against a fixed schema; invalid (malformed) outputs are retried deterministically until a valid integer is obtained ($\leq 3$ retries). These tokens are the basis for uncertainty feature computations~\cite{kapoor2024large,farr-etal-2025-red,gupta2024language,orgad2025llmsknowshowintrinsic}.\footnote{Prior work has also proposed an alternative strategy of adding a dedicated uncertainty token (e.g., ``[IDK]''), enabling the model to explicitly allocate probability mass to ``I don’t know'' rather than forcing a distribution over fixed labels~\cite{cohen2024i}.} Reasoning tokens are handled separately: they remain free-form and are used only for deriving complementary reasoning-based features (e.g., perplexity, per-token entropy), not for outcome probability computation~\cite{guerreiro-etal-2023-looking,efficient-and-effective}; (3) \firstmention{LPP Feature Extraction}: a comprehensive set of LLM Performance Predictors (LPPs) is computed from model outputs, leveraging {\em gray-box} access (token-level log-probabilities, entropy, and reasoning-path statistics) together with {\em black-box} compatible features (Verbalized Confidence and Uncertainty Attribution Indicators; see Section~\ref{subsec:lpps}); (4a) \firstmention{Meta-Model Training}: a Ridge Regression classifier is trained on LPPs to predict LLM correctness, providing calibrated uncertainty estimates~\cite{arthur+al:00}; and (4b) \firstmention{Cost-Aware Routing}: a threshold-based policy decides whether to trust the LLM prediction or escalate to human review, optimizing operational costs.
This architecture is inspired by Query Performance Prediction (QPP) in Information Retrieval~\cite{carmel+Elad:10}, adapted to the agent-based setting where a meta-classifier acts as a gating agent, coordinating between an LLM agent and human reviewers.
% requires: \usetikzlibrary{positioning}
\begin{figure*}[t]
\centering
\begin{tikzpicture}[
  font=\footnotesize,
  node distance=10mm and 8mm, % layout spacing only (not size)
  box/.style={draw, rounded corners, semithick, align=center, fill=#1!10},
  arr/.style={-Latex, semithick}
]

% Stage 1
\node[box=blue] (s1) {%
  \textbf{Stage 1}\\
  Base LLM inference\\
  \scriptsize multimodal input\\
  \scriptsize structured prompts
};

% Stage 2
\node[box=green, right=of s1] (s2) {%
  \textbf{Stage 2}\\
  Integer-token output schema\\
  \scriptsize \texttt{0,1,2,3} labels
};

% Stage 3
% \node[box=yellow, right=of s2] (s3) {%
%   \textbf{Stage 3}\\
%   LPP feature extraction\\
%    $\bullet$ Outcome probs\\
%    $\bullet$ Verbalized conf.\\
%    $\bullet$ Abstention flags
% };
% Stage 3
\node[box=yellow, right=of s2, text width=0.14\textwidth] (s3) {%
  \centering\textbf{Stage 3}\\
  \centering LPP feature extraction\\[-0.25ex]
  \raggedright\scriptsize
  \begin{tabular}{@{}l@{}}
    \textbullet\ Outcome probs\\
    \textbullet\ Verbalized conf.\\
    \textbullet\ Abstention flags
  \end{tabular}
};


% Stage 4a
\node[box=orange, right=of s3] (s4a) {%
  \textbf{Stage 4a}\\
  Meta-model training\\
  \scriptsize Ridge; outputs $s_\theta(x)$
};

% Stage 4b
\node[box=red, right=of s4a] (s4b) {%
  \textbf{Stage 4b}\\
  Cost-aware routing\\
  \scriptsize $s_\theta(x)\!\ge\!\tau$\\
  \scriptsize \emph{trust} vs. \emph{escalate}
};

% Arrows
\draw[arr] (s1) -- (s2);
\draw[arr] (s2) -- (s3);
\draw[arr] (s3) -- (s4a);
\draw[arr] (s4a) -- (s4b);

\end{tikzpicture}

\caption{A four-stage framework:
\textbf{(1)} Base LLM Inference;
\textbf{(2)} Integer-Token Output Schema;
\textbf{(3)} LPP Feature Extraction;
\textbf{(4a)} Meta-Model Training to estimate $s_\theta(x)$;
\textbf{(4b)} Cost-Aware Routing to \emph{trust} or \emph{escalate}.}
\label{fig:pipeline}
\end{figure*}
\subsection{LLM Performance Predictors (LPPs)}
\label{subsec:lpps}
Our LLM Performance Predictors (LPPs) are a comprehensive feature set primarily extracted via \textbf{gray-box access}—requiring token-level log-probabilities and structured outputs. This feature set also incorporates black-box compatible features (Verbalized Confidence and Uncertainty Attribution Indicators) derived solely from the structured text output. This strategy ensures broad applicability: all evaluated base LLMs provide the necessary gray-box and black-box signals. The combination of gray-box and black-box features strikes a practical balance between information richness and scalability, avoiding reliance on internal activations (white-box) while providing richer signals than output-only methods. Recent analysis has shown that probabilistic confidence (derived from token probabilities) and verbalized confidence capture complementary aspects of model uncertainty, with the former being more accurate but requiring threshold calibration, while the latter provides reasonable signal without additional setup~\cite{ni+al:25}.
LPPs span four families, summarized in Table~\ref{tab:lpps}, and the full mathematical specification of all LPPs appears in Table~\ref{app:full_lpps} (Appendix~\ref{app:prompts}). 
\myparagraph{Post-Hoc Classification Uncertainty}
Features derived from the probability distribution over outcome tokens, computed over the top-$k$ most probable tokens ($k=5$), including Top-5 Entropy $H(\tilde{p})$, Normalized Top-5 Entropy, Effective Choices ($2^{H(\tilde{p})}$), Max Softmax Probability, and Top-2
Probability Margin~\cite{kapoor2024large,kendall+al:17,farr+al:25,watson2023explaining,shorinwa2025surveyuncertaintyquantificationlarge, correa+al:25}.
\myparagraph{Internal Generation Uncertainty}
We implemented Chain-of-Thought (CoT) prompting solely to extract reasoning-sequence features (e.g., Perplexity, Sequence Negative Log-Likelihood, Mean Token Entropy, Token Entropy Quantiles, and Token Probability Quantiles)~\cite{guerreiro-etal-2023-looking,efficient-and-effective,correa+al:25}. 
In our content-moderation setting, CoT inflated confidence and harmed calibration; therefore, these features are documented for completeness but excluded from reported results.
\myparagraph{Verbalized Confidence}
Self-reported confidence extracted from structured outputs: a scalar reported confidence $\hat{c} \in [0,100]$ (normalized to $[0,1]$) and coarse confidence bands (Very Low, Low, Medium, High, Very High) encoded as one-hot features~\cite{kadavath+al:22,tian+al:23,yang2024verbalizedconfidencescoresllms,shorinwa2025surveyuncertaintyquantificationlarge,ni+al:25,xiong+al:23}.
\myparagraph{Uncertainty Attribution Indicators}
Binary indicators when the LLM emits \textsf{inconclusive\_evidence} (aleatoric uncertainty due to missing context, ambiguous frames, etc.) or \textsf{inconclusive\_definition} (epistemic uncertainty due to policy gaps or edge cases). These novel moderation-oriented features enable interpretable routing~\cite{kendall+al:17}: evidence deficient cases escalate to senior reviewers, while policy ambiguous cases trigger guideline updates or active learning cycles.
\begin{table}[h]
\centering
\caption{\textbf{Overview of LPP Feature Families.} Each category captures a distinct aspect of model uncertainty and is accessible via gray-box or black-box access. The full mathematical specification of all LPPs appears in Table~\ref{app:full_lpps} (Appendix~\ref{app:prompts}).}
\small
\setlength{\tabcolsep}{3.5pt}
\renewcommand{\arraystretch}{1.05}
\begin{tabularx}{\columnwidth}{@{}%
  >{\raggedright\arraybackslash}p{0.24\columnwidth}%
  >{\raggedright\arraybackslash}X%
  >{\raggedright\arraybackslash}p{0.3\columnwidth}%
@{}}
\toprule
\textbf{LPP Category} & \textbf{Signal Source \& Intuition} & \textbf{Example Features} \\
\midrule
\textbf{Post-Hoc Classification Uncertainty} (Gray-Box) &
Confidence at the final decision boundary, derived from token log-probabilities over discrete classification outcomes. &
Max Softmax Probability (MSP), Top-5 Entropy (Entropy), Top-2 Probability Margin (Top-2 Margin), Confidence Score \\
\midrule
\textbf{Internal Generation Uncertainty} (Gray-Box; excluded from reported results) &
Signals from intermediate Chain-of-Thought (CoT) reasoning tokens. &
Perplexity (Natural Base), Sequence Negative Log-Likelihood, Mean Token Entropy, Token Entropy Quantiles \\
\midrule
\textbf{Verbalized Confidence} (Black-Box) &
Explicit, self-reported confidence stated by the LLM in structured natural language outputs. &
Reported Confidence (Scalar), Confidence Bands (One-Hot; VL/L/M/H/VH) \\
\midrule
\textbf{Uncertainty Attribution Indicators} (Black-Box) &
Uncertainty attribution indicators signals distinguishing aleatoric uncertainty (evidence deficits) from epistemic uncertainty (policy gaps). Enables interpretable routing. &
Inconclusive Flag (Binary), Evidence-Deficit Indicator, Policy-Gap Indicator \\
\bottomrule
\end{tabularx}
\label{tab:lpps}
\end{table}
\vspace{-11pt}

\subsection{Prompting}
\label{subsec:prompting}
We designed four prompt templates: text-only and multimodal, each with two variants, a direct-answer version (no explicit reasoning steps) and a Chain-of-Thought (CoT) version. All prompts required a structured JSON output with (i) a single integer label (\texttt{0--3}), (ii) an optional reasoning field (populated only in the CoT variant), and (iii) any self-reported confidence. The integer-coded schema aligns labels with token-level log-probabilities. For completeness, we implemented both variants; unless otherwise stated, all results reported in Section~\ref{sec:results} use the direct-answer variant.
\subsection{Cost-Aware Escalation Policy}
\label{subsec:policy}
The meta-model output is a score $s_\theta(x) \in [0,1]$ representing the probability that the base LLM is correct. To operationalize this into a binary trust-or-escalate decision, we employ a \textbf{cost-aware threshold policy}. A prediction is trusted if $s_\theta(x) \geq \tau$; otherwise, it is escalated to human review~\cite{elkan+al:01}.
Let $c_{\text{mis}}$ denote the cost of an undetected misclassification (allowing an LLM error into production) and $c_{\text{rev}}$ the cost of human review. The expected cost, relative to a baseline of always trusting the LLM, is:
\[
\mathcal{C}(\tau) = c_{\text{mis}} \cdot \text{FP} + (c_{\text{rev}} - c_{\text{mis}}) \cdot \text{TN} + c_{\text{rev}} \cdot \text{FN},
\]
where FP, TN, FN are counts of selector decisions (trust/escalate) against LLM correctness (correct/incorrect). Here:
\begin{itemize}
    \item \textbf{TP (True Positives)}: Trusting a correct LLM prediction $\Rightarrow$ no costs.
    \item \textbf{FP (False Positives)}: Trusting an incorrect LLM prediction $\Rightarrow$ misclassification cost (error is missed).
    \item \textbf{TN (True Negatives)}: Escalating an incorrect LLM prediction $\Rightarrow$ review cost offset by avoided misclassification cost (error is caught).
    \item \textbf{FN (False Negatives)}: Escalating a correct LLM prediction $\Rightarrow$ unnecessary review cost.
\end{itemize}
The cost of human reviewers was calculated by multiplying the total review time by their hourly rate (in \$). The cost of misclassification was estimated based on projected business losses, specifically the reduction in customer lifetime value caused by churn due to undetected errors. In our experiments, we set the cost ratio such that $c_{\text{rev}} / c_{\text{mis}} \approx 0.64$, which serves as a reasonable baseline and does not alter relative rankings or cost–accuracy trends under variation.
The operating threshold $\tau^{\star}$ is optimized on a held-out validation set (20\% of the training data) by sweeping $\tau \in [0.35, 0.70]$, a practical range within which the optimum consistently lies near the midpoint (0.5). The value minimizing $\mathcal{C}(\tau)$ is then fixed and applied to the held-out test set for evaluation.
\subsection{Meta-Model Training}
\label{subsec:training}
Our meta-model is a supervised binary classifier trained to predict whether the base LLM is correct ($z=1$) or incorrect ($z=0$) given the extracted LPPs. Based on extensive preliminary experiments, we selected \textbf{Ridge Regression}, whose output scores were converted into probabilities using either sigmoid or isotonic calibration via scikit-learn's \texttt{CalibratedClassifierCV}. This approach was chosen for the following reasons:
\myparagraph{Calibration and Robustness} Ridge Regression's $L_2$ regularization provides crucial protection against multicollinearity, expected within the LPP feature set where signals like entropy and max probability are inherently correlated. This regularization prevents overfitting and improves out-of-sample calibration, critical for reliable uncertainty quantification in production systems.
\myparagraph{Interpretability} Ridge coefficients can be inspected to understand which LPPs are most predictive of errors, supporting explainability requirements in high-stakes moderation workflows.
\myparagraph{Gray-Box Applicability} Ridge Regression does not require model internals (hidden states, gradients), ensuring compatibility with both off-the-shelf APIs and open-source LLMs.
\myparagraph{Handling Class Imbalance}
In the Multimodal Moderation Dataset (1,500 samples), base LLM predictions were correct in roughly 80--90\% of cases (10--20\% errors), resulting in a strong class imbalance. To address this, we (1) downsampled the majority class. We applied a hybrid undersampling strategy combining Tomek Links~\cite{wonjae+al:22} for boundary cleaning with random undersampling of the remaining majority samples. This reduces redundancy while maintaining representative class ratios, and ensures that rare abstention cases ($<5\%$ of data) are always retained. We also tested oversampling methods such as SMOTE~\cite{chawla+al:02}, which yielded comparable performance; for simplicity and reproducibility, we report results using the Tomek+random undersampling approach. We also tuned Ridge’s \texttt{class\_weight} via grid search, testing ratios informed by the cost structure: $w_0 / w_1 \approx c_{\text{rev}} / c_{\text{mis}} \approx 0.64$.
After downsampling and stratified train–test splitting, our typical training set contains $\sim$800 samples and the test set $\sim$300 samples for the Multimodal Moderation Dataset, and $\sim$3{,}500 training and $\sim$900 test samples for the OpenAI Moderation Dataset, with both classes well represented. To ensure fair comparison across LLMs under class imbalance, we fixed the number of negative (error) cases in the test sets: 45 for the Multimodal Moderation Dataset and 150 for the OpenAI Moderation Dataset. This controlled evaluation design mitigates variability due to scarcity of negative samples while preserving comparability across models.
\myparagraph{Cross-Validation and Hyperparameter Search}
We employ stratified $3$-fold cross-validation and a comprehensive grid search on the training portion of the data (after holding out a validation split for threshold selection) over:
\begin{itemize}
    \item \texttt{alpha} $\in \{0.1, 1.0, 10.0, 100.0\}$ (regularization strength)
    \item \texttt{solver} $\in \{\text{auto}, \text{lsqr}\}$
    \item \texttt{tol} $\in \{10^{-6}, 10^{-5}, 10^{-4}, 10^{-3}\}$ (convergence tolerance)
    \item \texttt{max\_iter} $\in \{1000, 2000, 3000\}$
    \item \texttt{class\_weight}: 7 configurations including cost-informed ratios and ``balanced''
    \item Calibration \texttt{method} $\in \{\text{sigmoid, isotonic}\}$
\end{itemize}
Models are selected based on F1-score on the minority class (errors), as this metric balances precision and recall for the operationally critical decision of identifying when the LLM is wrong.\\
Formally, each sample is labeled with a binary correctness indicator $z_i \in \{0,1\}$, equal to $1$ if the LLM prediction matches the ground truth and $0$ otherwise.  
The ridge regression objective is
\[
\min_{w,b} \; \frac{1}{N}\sum_{i=1}^N \big(z_i - (w^\top f(x_i) + b)\big)^2 + \lambda \|w\|_2^2,
\]
where $f(x_i) \in \mathbb{R}^d$ is the feature vector, $w \in \mathbb{R}^d$ and $b \in \mathbb{R}$ are model parameters, and $\lambda \ge 0$ controls $L_2$ regularization. This yields a calibrated linear predictor of correctness probability, with isotonic or Platt scaling to map scores into $[0,1]$.
Inference settings (including deterministic decoding) are described in Section~\ref{subsec:experimental_setup} and Appendix~\ref{app:prompts}.